\epigraph{To doubt everything or to believe everything are two equally convenient solutions; both dispense with the necessity of reflection.}
{\textit{Henri Poincaré}}

\section{First Year Mathematics Workshops}

\subsection{Introduction to Workshops}

%okay - check with Armando on what it the most accurate up to date info on this - the webiste MATH1046 is not really updated 

% Talk formally about how often they run - how they run - who's present to help, etc. any other useful info. 

% also stress how lucky we are as a dept to have these things, some unis don't give a fuck about students. 



\subsection{More Than Just Coursework Help}

When I first started, I thought the first year workshops were basically just getting help for your coursework, and nothing else. 

But that would be missing a major point. Workshops are not there to replace your thinking, but to refine it. 

The biggest difference I noticed was that when I genuinely struggled with a problem, even if I got nowhere, the workshop became the most valuable resource. If you go hoping someone will just show you the solution from scratch, you leave having learned much less than you could have. Even 15 minutes of serious attempt make major impact. It makes you aware of where you're stuck, and gives you something specific to ask about. 

\section{AI and Outsourcing Thinking}

It would be highly ignorant and unrealistic of me to not talk about Generative AI. We all know it exists and it can generate solutions instantly. And sometimes, it is useful, it can clarify notation or rephrase definitions and help fill in some gaps in reasoning.

However, the biggest drawback of AI is that it removes the part of the process that actually builds the understanding. If you ask AI to generate a full solution to a problem you haven't seriously attempted, you might feel productive, even pleased looking at the solution, but you've just skipped the uncomfortable but crucial bit where your brain learns by trying, failing, and reattempting ideas. And that's the uncomfortable bit where understanding takes place. 

There is a difference between using AI to complement your learning, and using it to replace your work and your thinking entirely, and it is important that you be honest with yourself about which one you're doing. At the end of the day, you are the one being impacted by your use/misuse of AI, so take accountability. 

The department and the university has its own policy on AI, so I urge you to read it because it matters not only due to academic integrity, but because this degree is about learning to think independently. 
% There are univeristy and department policies on AI, add links to them so that I don't end up saying something that will be hugley incorrect? 


\section{Beyond Workshops}

% peers, staff office hours, tutor, older students, etc. 


